\clearpage\begingroup\SourceLeaf{407b}\ArithmeticHead{76}{L}{STÄCKEL, GAUSS ALS GEOMETER.}
\renewcommand{\thefootnote}{\arabic{footnote})}\setcounter{footnote}{2}
\begin{center}
\textbf{25.}\par\smallskip
\textls[100]{Der Kreis und die Kugel.}
\end{center}
\noindent\hspace*{\parindent}Das \textit{Tagebuch} von \textsc{Gauss} beginnt mit der Eintragung vom 30. März 1796:
»Principia quibus innititur sectio circuli, ac divisibilitas eiusdem geometrica
in septemdecim partes etc«. Nach dem Briefe an \textsc{Gerling} vom 6. Januar 1819
hatte er die Entdeckung am Morgen des 29. März 1796 gemacht (W. X\textsubscript{1},
S. 125). »Sie ist es vornehmlich gewesen, welche seinem Leben eine bestimmte
Richtung gab, denn von jenem Tage an war er fest entschlossen,
nur der Mathematik sein Leben zu widmen« (\textsc{Sartorius}, S. 16). Die Konstruktion
des regelmässigen Siebzehnecks ist geometrisch ausführbar, insofern
sie sich allein durch Lineal und Zirkel bewerkstelligen lässt, jedoch beruht
der Beweis bei all’ den verschiedenen Durchführungen auf der algebraischen
Grundlage der Kreisteilungsgleichung\footnote{Vgl. R. \textsc{Goldenring}, \textit{Die elementargeometrischen Konstruktionen des regelmässigen Siebzehnecks}, Dissertation, Jena 1915. Die zeitlich älteste Konstruktion ist die dort noch nicht erwähnte von \textsc{Pfleiderer}, die erst 1917 (Werke X\textsubscript{1}, S. 120) veröffentlicht worden ist.}. \textsc{Gauss} hat in den Göttinger Anzeigen
vom 19. Dezember 1825 eine Konstruktion von \textsc{Erchinger} mitgeteilt. Für
diese Konstruktion habe \textsc{Erchinger} eine rein geometrische Begründung gegeben,
»mit musterhafter, mühsamer Sorgfalt, alles nicht rein Elementarische
zu vermeiden« (W. II, S. 187). Sie ist uns leider verloren gegangen, da \textsc{Erchingers}
Abhandlung nicht gedruckt wurde\footnote{Die Abhandlung \textsc{Erchingers} hatte \textsc{Gauss} von einem Braunschweiger Bekannten, dem Juristen
E. \textsc{Schrader} in Tübingen, am 1. Sept. 1825 zugesandt erhalten. Hiernach war \textsc{Erchinger}, der sonst ganz
unbekannt ist, ein mathematischer Autodidakt, der etwa seit 1813 in Tübingen lebte. Er hatte einen Beitrag
geliefert zu der Abhandlung \textsc{Schraders}: \textit{Commentatio de summatione seriei}
\[
\frac{a}{b(b+d)}+\frac{a}{(b+2d)(b+3d)}+\frac{a}{(b+4d)(b+5d)}+\cdotp\,\cdotp\,\cdotp,
\]
Weimar 1818, die einen Preis der Kopenhagener Gesellschaft der Wissenschaften erhalten hatte. Nach
\textsc{Schraders} Brief an \textsc{Gauss} vom 20. April 1831 war \textsc{Erchinger} inzwischen gestorben (Briefe im \textsc{Gauss}-Archiv).
Vgl. auch \textsc{Klügels} \textit{Mathematisches Wörterbuch}, IV. Teil, Leipzig 1823, S. 652 (Artikel Summierung
der Reihen).}.
\clearpage\SourceLeaf{408}\ArithmeticHead{77}{R}{ELEMENTARE UND ANALYTISCHE GEOMETRIE.}\setcounter{footnote}{0}
\noindent\hspace*{\parindent}In dem Zusatz VI zu \textsc{Carnots} \textit{Geometrie der Stellung} (2. Teil, S. 371,
W. IV, S. 399) wird eine analytische Lösung der Aufgabe gegeben, einen
Kreis zu beschreiben, der drei der Grösse und Lage nach gegebene Kreise
berührt, »vielleicht die einfachste Konstruktion des Apollonischen Problems«,
wie \textsc{Simon} sagt\footnote{M. \textsc{Simon}, \textit{Über die Entwicklung der Elementargeometrie im XIX. Jahrhundert}, I. Ergänzungsband des Jahresberichtes der Deutschen Mathematiker-Vereinigung, Leipzig 1906, S. 98; man findet hier (S. 97—105) eine Zusammenstellung der umfangreichen Literatur über das Apollonische Taktionsproblem.}; sie ist wiederum der Benutzung komplexer Grössen zu verdanken.
\par
Um das Jahr 1840 hat \textsc{Gauss} den Begriff der harmonischen Punktepaare
auf einer Geraden verallgemeinert, indem er die vier Abszissen als komplexe
Grössen auffasst, denen Punkte einer Ebene zugeordnet sind (W. VIII, S. 336—337).
Hierin liegt ein fruchtbares Übertragungsprinzip, das \textsc{Möbius}, hier
wiederum mit \textsc{Gauss} zusammentreffend, ausgebaut hat\footnote{A. F. \textsc{Möbius}, \textit{Über eine Methode, um von Relationen, welche der Longimetrie angehören, zu den entsprechenden Sätzen der Planimetrie zu gelangen}, Leipziger Berichte, Bd. 4, 1852, S. 41, Werke II, S. 189.}. Später ist \textsc{Möbius}
zum allgemeinen Doppelverhältnis übergegangen und zu seiner Lehre von der
Kreisverwandtschaft gelangt, bei der zwischen zwei Ebenen durch eine bilineare
Gleichung in den lagebestimmenden komplexen Grössen eine Beziehung hergestellt
wird\footnote{A. F. \textsc{Möbius}, \textit{Über eine neue Verwandtschaft zwischen ebenen Figuren}, Leipziger Berichte, Bd. 5, 1853, S. 14, Werke II, S. 205; später hat \textsc{Möbius} die Kreisverwandtschaft rein geometrisch begründet: \textit{Die Theorie der Kreisverwandtschaft in rein geometrischer Darstellung}, Leipziger Abhandlungen, Bd. 4, 1855, S. 529; Werke II, S. 243.}.
\par
Auch den Punkten einer Kugelfläche hat \textsc{Gauss} schon sehr früh komplexe
Grössen zugeordnet, vermutlich mittels der stereographischen Projektion; dies
zeigt die schon angeführte Bemerkung über das elliptische Integral erster
Gattung aus dem Jahre 1800 (W. X\textsubscript{1}, S. 546). In einer späteren
\clearpage\SourceLeaf{409a}\ArithmeticHead{78}{L}{STÄCKEL, GAUSS ALS GEOMETER.}\setcounter{footnote}{0}
\noindent Aufzeichnung, die vor 1819 niedergeschrieben ist, hat er den durch die stereographische
Projektion vermittelten Zusammenhang zwischen Ebene und Kugel genauer
untersucht und dabei erkannt, dass die »Drehungen der Kugelfläche in sich
selbst« durch gewisse lineare, gebrochene Substitutionen der lagebestimmenden
komplexen Grösse dargestellt werden können (W. VIII, S. 854—356); man
kennt die Bedeutung, die diese Substitutionen später gewonnen haben\footnote{Vgl. für die von \textsc{Riemann} benutzte Verwendung der Kugel zur Darstellung komplexer Grössen C. \textsc{Neumann}, \textit{Vorlesungen über Riemanns Theorie der Abelschen Integrale}, Leipzig 1865, für die linearen Substitutionen F. \textsc{Klein}, \textit{Vorlesungen über das Ikosaeder}, Leipzig 1884, erster Abschnitt, Kapitel II.}.
\clearpage\endgroup
